\subsection{An ordinary new factor with an odd relative class number}
\label{sec:ordinary-rational-witness}

\paragraph{Finding A32 (the rational witness can be geometrically ordinary).}
Let $C/\mathbf F_4$ be the curve of A27 and A31, let $W$ be its
point over $t=0$, and put $K=\mathbf F_{256}$.  There is a constant,
geometrically connected, $W$-split cyclic etale cover $D\to C_K$
of degree thirteen whose Prym
\[
             \mathcal P=\ker(J_D\longrightarrow J_C)^0
 \tag{OR.1}
\]
is ordinary of dimension twelve, while $\#\mathcal P(K)$ is odd.
Writing $K_m=\mathbf F_{256^m}$, the exact extension rule is
\[
 \frac{\#J_D(K_m)}{\#J_C(K_m)}\ \text{is odd}
             \quad\Longleftrightarrow\quad 7\nmid m.
 \tag{OR.2}
\]
All $2^{12}$ geometric etale $2$-torsion points of $\mathcal P$
are rational over $K_7$; seven is the least extension degree
with any nonzero rational $2$-torsion in that factor.
The reduced Frobenius numerator is
\[
                H_{\mathcal P}(T)\bmod2=1+T^4+T^{12}.
 \tag{OR.3}
\]
Thus this rational class-number witness uses a geometrically
ordinary new factor.  It is a different mechanism from A31's
geometric rank-zero witness, and it applies at a Weierstrass
point where the canonical class detector of A29 is zero.

\begin{proof}
\emph{1. The exact cyclic quotient and its point-image table.}
For the calculation use the isomorphic model
\[
 C_0:\quad Y^2+X^2Y=uX^5+X^3+uX,
 \qquad X=1/t,\quad Y=y/t^2.
 \tag{OR.4}
\]
Its point at infinity is $W$.  Realize
\[
 K=\mathbf F_2[z]/(z^8+z^4+z^3+z+1),\qquad
 u=z^7+z^5+z^4+z^3+z^2.
 \tag{OR.5}
\]
Then $u^2+u+1=0$.  For the finite table only, an integer
$\sum_i a_i2^i$ encodes the field element $\sum_i a_iz^i$.
Thus $u$ is encoded by $188$.

Finding A27 gives
\[
        G=J_C(K),\qquad |G|=65520=13\cdot5040,
        \qquad |C(K)|=256.
 \tag{OR.6}
\]
The thirteen-primary subgroup is therefore cyclic of order
thirteen.  In Mumford coordinates for (OR.4), the element
\[
        \Delta=(X^2+80X+237,\ 106X+207)
 \tag{OR.7}
\]
has exact order thirteen; the coefficients in this display use
the field encoding just specified.  Define the surjection
$a:G\to\mathbf Z/13$ by
\[
                      [5040]Q=a(Q)\Delta.
 \tag{OR.8}
\]
Multiplication by $5040$ kills every other primary factor and
is an automorphism on the thirteen-primary subgroup.  Exact
enumeration gives the following complete distribution:
\[
\begin{array}{c|rrrrrrrrrrrrr}
 j&0&1&2&3&4&5&6&7&8&9&10&11&12\\ \hline
 \#\{R:a([R-W])=j\}
 &44&19&18&18&16&19&16&16&19&16&18&18&19.
\end{array}
 \tag{OR.9}
\]
The entries add to $256$, including the point at infinity.

Here the finite computation has two separate implementations.
The first uses the extended-GCD composition formula for Mumford
divisors.  The second imports none of its arithmetic and obtains
composition by binary linear algebra as follows.  Write a reduced
class as $(U,V)$, where $U$ is monic, $\deg U\le2$,
$\deg V<\deg U$, and
\[
                         U\mid V^2+X^2V+f,
                \qquad f=uX^5+X^3+uX.
 \tag{OR.10}
\]
To add two such classes, cancel their common opposite support
using $\gcd(U_1,U_2,V_1+V_2+X^2)$.  This removes pairs whose
sum is the divisor of their $X$-polynomial, hence does not change
the resulting Jacobian class.  More explicitly, a nonramified
pair $Q+\iota Q-2W$ is the divisor of $X-X(Q)$, with the
same cancellation repeated for each common multiplicity.  The
only affine branch point is $(0,0)$; it occurs at most once in
a reduced pair, since $f'(0)=u\ne0$ prevents $X^2$ from
dividing $V^2+X^2V+f$.  A double occurrence therefore cancels
by $\operatorname{div}(X)=2(0,0)-2W$.
For the remaining polynomials,
solve for $V$ of degree below $\deg(U_1U_2)$ subject to
\[
 \begin{aligned}
 V&\equiv V_i\pmod{U_i}\quad(i=1,2),\\
 V^2+X^2V+f&\equiv0\pmod{U_1U_2}.
 \end{aligned}
 \tag{OR.11}
\]
This is a linear system over $\mathbf F_2$ in the binary
coefficients of $V$.  Common remaining support is on the same
unramified branch, where the derivative in $Y$ is nonzero;
the required higher-multiplicity root interpolation is therefore
unique, by recursively solving the local root equation using
that nonzero derivative.  This applies also to a doubled
degree-two divisor; it is not restricted to distinct support.
At the ramified point any opposite pair was already cancelled,
and a remaining simple occurrence fixes only its value.
Consequently the resulting pair represents the sum divisor
with its correct multiplicities.

If its degree is above two, replace it by
\[
     U'=(V^2+X^2V+f)/U\quad\text{made monic},\qquad
     V'=V+X^2\pmod{U'}.
 \tag{OR.12}
\]
Here the divisor of $Y+V$ has finite part $D_U+E$ and pole
order $\max\{5,2\deg V\}$ at $W$.  The two possible pole
orders are odd and even, so there is no leading cancellation.
Its norm is $U$ times the polynomial for $E$.  Thus
$[D_U-(\deg U)W]=-[E-(\deg E)W]$, and replacing $E$ by its
hyperelliptic conjugate gives precisely the second formula in
(OR.12).  The leading coefficient $u$ of $f$ is retained in
the division; making $U'$ monic multiplies only its defining
polynomial by a nonzero constant and changes no divisor.
The new degree is at most two.  A reduced class is unique:
pad its effective representative to degree two with $W$ and
apply genus-two Riemann--Roch.  Two distinct representatives
could occur only in the canonical degree-two pencil; its
conjugate pairs have already been cancelled to the identity.
The second implementation verifies (OR.10) after every step,
uses a different scalar-multiplication chain, enumerates all
affine solutions of (OR.4) plus infinity, and independently
reproduces (OR.7)--(OR.9).  Its companion point file retains all
$256$ coordinates and quotient images.  Thus the finite input
to the character calculation is fully specified and reproducible.

\emph{2. The pointed cover and the reduced character values.}
Take the connected constant Lang quotient cover associated to
(OR.8), based at $W$, as in A12.  Its fibre over $W$ is split.
The genus is fourteen, so the dimension of $\mathcal P$ is twelve.
Since $K_C\sim2W$, the exact genus-two polynomial of recovered
A1 B.2 is
\[
 L_W(\chi,T)=1+c_1(\chi)T+256T^2,
 \qquad
 c_1(\chi)=\sum_{R\in C(K)}\chi(a([R-W]))
 \tag{OR.13}
\]
for every nontrivial character of $\mathbf Z/13$
\cite[B.2]{A1run}.  Put $\zeta$ for a primitive thirteenth root
in characteristic two.  Formula (OR.9) gives
\[
 \begin{aligned}
 S(Z)&=Z+Z^5+Z^8+Z^{12},\\
 \overline{L_W(\chi_i,T)}&=1+S(\zeta^i)T
                         \qquad(1\le i\le12).
 \end{aligned}
 \tag{OR.14}
\]
In particular the B.2 group-ring test is $1+S(Z)$.
Two has order twelve modulo thirteen, so
$\Phi_{13}(Z)=1+Z+\cdots+Z^{12}$ is irreducible over
$\mathbf F_2$.  Neither $S$ nor $1+S$ is that polynomial or
zero.  Therefore
\[
                        S(\zeta^i)\notin\{0,1\}
                           \quad(1\le i\le12).
 \tag{OR.15}
\]
The simultaneous test of B.2 now gives an odd relative class
number.  The norm-pullback isogeny identifies that integer with
$\#\mathcal P(K)$, as in B.1.

\emph{3. Exact character scalars and ordinary geometric rank.}
Set $\beta=S(\zeta)$.  Multiplication of the four exponents
$\{1,5,8,12\}$ by eight permutes them modulo thirteen, so
$\beta^8=\beta$.  Thus $\beta\in\mathbf F_8$, and (OR.15)
shows that it has exact multiplicative order seven.
The three cosets obtained by multiplying those exponents by
$1,2,4$ partition $(\mathbf Z/13)^\times$.  It follows that
\[
              \beta+\beta^2+\beta^4
                        =\sum_{i=1}^{12}\zeta^i=1.
 \tag{OR.16}
\]
Its minimal polynomial is therefore the trace-one irreducible
cubic $Z^3+Z^2+1$.  Consequently
\[
 \prod_{j=0}^{2}(1+\beta^{2^j}T)=1+T+T^3,
 \qquad
 \prod_{i=1}^{12}(1+S(\zeta^i)T)=(1+T+T^3)^4.
 \tag{OR.17}
\]
The twelve character scalars comprise the three conjugates,
each repeated four times.  Character factorization in B.2
now proves (OR.3) directly.

For clarity, the cohomological interpretation is a cited result,
not a new comparison theorem.  Recovered successor A.20 identifies
the reduced character polynomial with the determinant of the
linear extension of $256$-power coherent Frobenius on the
character summand of $H^1(D,\mathcal O_D)$
\cite[A.20]{ASrun}.  Its etale determinant step is Deligne,
SGA~$4\frac12$, \emph{Fonctions $L$ modulo $\ell^n$ et modulo
$p$}, Theorem~2.2(b) and formula~(2.4.1), with the hypotheses
of \S2.1 \cite{DeligneModP}.  Here the coefficient ring is a
finite field of characteristic two containing the character
values, hence is noetherian and reduced; the rank-one lisse
character sheaf is constructible and has finite Tor dimension
on the smooth proper curve $C_K$.  Each nontrivial summand has dimension one
because the base genus is two.  All its displayed scalars are
nonzero, so the full twelve-dimensional new part is stable under
Frobenius.  The Artin--Schreier stable-part comparison therefore
gives geometric $2$-rank twelve; equivalently, the new factor is
ordinary.  The exact primary comparison locator is Katz,
SGA~7~II, Expos\'e~XXII, Proposition~2.2.5, with the
characteristic-polynomial convention in Proposition~2.2.4
\cite{KatzCongruence}.

\emph{4. Constant extensions and rational torsion.}
After extension to $K_m$, the linear Frobenius on those
one-dimensional summands has scalars $S(\zeta^i)^m$.
Every one of the original scalars has exact order seven.
Thus the relative class-number reduction is
\[
        \prod_{i=1}^{12}\bigl(1+S(\zeta^i)^m\bigr),
 \tag{OR.18}
\]
which is nonzero exactly when $7\nmid m$.  This proves (OR.2).
The deck group is constant, so Frobenius preserves the twelve
character lines; there is no extra Jordan block hidden by their
repeated characteristic-polynomial roots.  At $m=7$ its action
on the entire new stable part is the identity.  The functorial
stable-part comparison, or its dual on etale Jacobian torsion,
then makes all $2^{12}$ geometric etale $2$-torsion points of
$\mathcal P$ rational.  At smaller positive degrees none is
rational, by (OR.18).  This proves the claimed exact first
extension degree.
\end{proof}

At this same level the numerical premises are very far from
necessary: $|G/13G|=13<256=|C(K)|$, and B.3's pair bound would
require $256\ge2\cdot5040$.  Finding A32 supplies an actual
primitive-two prime witness despite both failures.  Its ordinary
new factor also shows why rational class-number units must not
be equated with geometric rank zero.  No absence of other
geometric theta points of order thirteen is asserted.
